% =====================================================================
%  2bat-ud2-6.tex  —  SESSIÓ 6: GeoGebra: dibuixar regions factibles
% =====================================================================
\horatitol{Sessió 6 — GeoGebra: dibuixar regions factibles}%
{Representar amb l'ordinador les regions factibles fetes a mà, comparar-les i localitzar els vèrtexs amb precisió.}

\subsection{Idea clau}

Avui portem a \textbf{GeoGebra} les regions factibles que hem dibuixat a mà.
L'aplicació ombreja sola el semiplà de cada inequació i, en posar-les totes,
mostra clarament la \textbf{regió factible}. La idea important: GeoGebra
\textbf{no fa la feina per tu}; serveix per \textbf{comprovar} la teva i per trobar
els \textbf{vèrtexs} amb precisió.

\begin{idea}
\textbf{Com representar un sistema a GeoGebra}
\begin{enumerate}[leftmargin=1.6em, itemsep=2pt]
  \item \textbf{Introdueix cada inequació} a la barra d'entrada, tal com s'escriu:
        \texttt{x+2y<=12}, \texttt{x+3y<=15}, \texttt{x>=0}, \texttt{y>=0}.
        GeoGebra ombreja el semiplà de cada una.
  \item \textbf{La regió factible} és la zona on se superposen \emph{tots} els
        ombrejats (la més fosca).
  \item \textbf{Vèrtexs:} usa l'eina \emph{Intersecció} (o l'ordre
        \texttt{Interseca}) sobre dues vores per obtenir-ne les coordenades exactes.
  \item \textbf{Compara} la regió de la màquina amb la que vas fer a mà: si no
        coincideixen, revisa el teu dibuix.
\end{enumerate}
\end{idea}

\subsection{Exemples}

\begin{exemple}[introduir un sistema i llegir-ne la regió]
Per al cas de l'entitat, escrivim a GeoGebra:
\[
\texttt{x+2y<=12}, \quad \texttt{x+3y<=15}, \quad \texttt{x>=0}, \quad \texttt{y>=0}.
\]
La zona on coincideixen tots els ombrejats és la regió factible. Amb l'eina
\emph{Intersecció}, en creuar \texttt{x+2y=12} amb \texttt{x+3y=15} obtenim el punt
$(6,3)$, exactament el vèrtex que havíem calculat a mà.

\begin{center}
\begin{tikzpicture}
\begin{axis}[udaxis, width=8.6cm, height=5.8cm,
    xlabel={\small $x$}, ylabel={\small $y$},
    xmin=0, xmax=15, ymin=0, ymax=8, xtick={0,4,8,12}, ytick={0,2,4,6}]
  \fill[udblue!18] (axis cs:0,0) -- (axis cs:12,0) -- (axis cs:6,3) -- (axis cs:0,5) -- cycle;
  \addplot[udblue, domain=0:12, samples=2]{(12-x)/2};
  \addplot[udnavy, domain=0:15, samples=2]{(15-x)/3};
  \addplot[udnavy, only marks, mark=*, mark size=1.5pt] coordinates {(6,3)};
  \node[font=\scriptsize, anchor=south west] at (axis cs:6,3) {$(6,3)$};
\end{axis}
\end{tikzpicture}

\figcap{La regió que mostra GeoGebra coincideix amb la feta a mà; la intersecció dóna el vèrtex $(6,3)$.}
\end{center}
\end{exemple}

\begin{exemple}[la màquina detecta un error]
Suposem que a mà havíem dibuixat la vora del material pel costat equivocat (com si
fos $x+3y\ge 15$). En posar-ho a GeoGebra, l'ombrejat de la màquina queda a l'altre
costat i la regió factible \emph{no} coincideix amb la nostra. Aquesta discrepància
ens avisa: havíem triat malament el semiplà. La comparació funciona com a
\textbf{autocorrecció}.
\end{exemple}

\subsection{Exercicis resolts}

\begin{exresolt}[interpretar el que mostra la pantalla]
En introduir $\;x+y\le 6,\;x\ge 0,\;y\ge 0$, GeoGebra ombreja un triangle. Quins
vèrtexs hauries d'obtenir amb l'eina Intersecció?
\solucio
Les vores són $x=0$, $y=0$ i $x+y=6$. Les interseccions de dues en dues donen:
$x=0$ amb $y=0$ $\rightarrow (0,0)$; $y=0$ amb $x+y=6$ $\rightarrow (6,0)$; $x=0$ amb
$x+y=6$ $\rightarrow (0,6)$.
\resposta{Vèrtexs $(0,0)$, $(6,0)$ i $(0,6)$.}
\end{exresolt}

\begin{exresolt}[comparar màquina i mà]
A mà has obtingut una regió amb vèrtexs $(0,0)$, $(12,0)$, $(6,3)$, $(0,5)$, però
GeoGebra te'n mostra una amb vèrtex $(0,6)$ en comptes de $(0,5)$. Què ha pogut
passar?
\solucio
El vèrtex superior de l'eix vertical el fixa la restricció \textbf{més exigent}. A
$(0,6)$, el material valdria $0+3\cdot 6=18>15$: no es compleix. El vèrtex correcte
és $(0,5)$ (on $x+3y=15$). Probablement a GeoGebra has oblidat d'introduir la
restricció del material, o l'has escrit malament; revisa-la.
\resposta{Falta (o està mal escrita) la restricció del material; el vèrtex bo és $(0,5)$.}
\end{exresolt}

\subsection{Exercicis per fer}

\noindent\textit{Full de pràctica tecnològica. Per a cada sistema: escriu el que vas
obtenir a mà, comprova-ho a GeoGebra i anota si coincideix o si has trobat un error.}

\begin{exercicis}
\begin{exlist}
  \item Introdueix a GeoGebra $\;x+y\le 5,\;x\ge 0,\;y\ge 0\;$ i comprova els tres
        vèrtexs amb l'eina Intersecció.
  \item Introdueix el sistema del cas de l'entitat ($x+2y\le 12$, $x+3y\le 15$,
        $x\ge 0$, $y\ge 0$) i comprova que els quatre vèrtexs coincideixen amb els que
        tens a mà.
  \item Representa $\;x\le 5,\;y\le 4,\;x\ge 0,\;y\ge 0\;$ i digues quina forma té la
        regió i quins són els seus quatre vèrtexs.
  \item Completa el full de validació amb tres sistemes a la teva elecció:
        \begin{center}
        \renewcommand{\arraystretch}{1.4}
        \begin{tabular}{p{3.6cm} p{3.4cm} p{2cm} p{2.6cm}}
        \toprule
        \textbf{Sistema} & \textbf{Vèrtexs a mà} & \textbf{Coincideix?} & \textbf{Error trobat} \\
        \midrule
         & & & \\[10pt]
         & & & \\[10pt]
         & & & \\
        \bottomrule
        \end{tabular}
        \end{center}
  \item \textbf{(Autocorrecció)} Si en algun sistema la regió de GeoGebra no coincideix
        amb la teva, escriu en una frase quin va ser l'error (semiplà equivocat?
        restricció oblidada? recta mal dibuixada?).
  \item \textbf{(Coherència)} Explica per què la regió factible és sempre la zona on se
        superposen \emph{tots} els ombrejats, i no la unió de tots ells.
\end{exlist}
\end{exercicis}

% ---------------------------------------------------------------------
\fullsolucions{Sessió 6}
\begin{solucions}
\begin{sollist}
  \item Vèrtexs $(0,0)$, $(5,0)$ i $(0,5)$; han de coincidir amb la versió manual.
  \item Vèrtexs $(0,0)$, $(12,0)$, $(6,3)$ i $(0,5)$; tots quatre han de coincidir.
  \item La regió és un \textbf{rectangle}; vèrtexs $(0,0)$, $(5,0)$, $(5,4)$ i $(0,4)$.
  \item Resposta oberta (full de validació personal). Es valora la coherència entre el
        càlcul a mà i la lectura de GeoGebra i l'anotació dels desajustos.
  \item Resposta oberta. Es valora identificar bé el tipus d'error i corregir-lo.
  \item Perquè cal complir \textbf{totes} les restriccions \emph{alhora}: un punt és
        factible només si és a tots els semiplans a la vegada, és a dir, a la
        \textbf{intersecció}. La unió inclouria punts que només compleixen una condició.
\end{sollist}
\end{solucions}
